\textbf{Proof.}
Work on the probability-one event supplied by thm:plugin-correspondence.  By prop:kappa-exhaustion, \(\mathcal B_\kappa\) is compact, and the preceding theorem establishes continuity of \(w\).  Let \(U\) be any neighborhood of the unique minimizer \(B^\star\).  The compact set \(\mathcal B_\kappa\setminus U\) contains no minimizer.  The extreme-value theorem therefore gives the derived gap
\[
 \gamma_U:=\min_{B\in\mathcal B_\kappa\setminus U}w(B)-v_{\mathrm{mm}}>0. \tag{1}
\]
Let \(\delta_m=\sup_B\lvert\widehat w(B)-w(B)\rvert\).  When \(2\delta_m<\gamma_U\), every \(B\notin U\) obeys
\[
 \widehat w(B)\ge w(B)-\delta_m
 \ge v_{\mathrm{mm}}+\gamma_U-\delta_m
 >v_{\mathrm{mm}}+\delta_m
 \ge\widehat w(B^\star).
\]
Thus \(\widehat{\mathcal A}_{\mathrm{mm}}\subset U\) eventually.  Since \(U\) was arbitrary, every measurable selection converges almost surely to \(B^\star\).

For the necessity of the singleton condition, take \(n=3\), \(f_i(w)=w_i\), and independent Bernoulli \((1/2)\) assignments.  Then
\[
 \Pi=\tfrac12I_6,
 \qquad \Omega=\tfrac14\begin{bmatrix}I_3&-I_3\\-I_3&I_3\end{bmatrix}.
\]
Let \(F_1=F_0=\delta_0\), \(H=\delta_{(0,0)}\), and \(\kappa=1/8\).  All pilot outcomes are zero and \(Q(0)=0\), so \(w(B)=\widehat w(B)=0\) for every feasible basis.  The two distinct projective points
\[
 B^{(2)}=B(\sqrt3e_1,\sqrt3e_2),
 \qquad B^{(3)}=B(\sqrt3e_1,\sqrt3e_3)
\]
are feasible because \(S_{B^{(j)}}=3I_2/4\) and \(\lambda_{\min}(3^{-1}S_{B^{(j)}})=1/4>\kappa\).  Selecting \(B^{(2)}\) at even pilot sizes and \(B^{(3)}\) at odd pilot sizes is measurable, satisfies the outer-set conclusion exactly, and has no point limit. \(\square\)